\subsubsection{一般方程组的化简}\label{subsec:general-system-reduction}

现在考虑一般的偏微分方程组, 在一个非特征曲面上指定初始数据. 将说明如何把这样的问题化为 \eqref{eq:ck-intro-quasilinear-system}, \eqref{eq:ck-intro-zero-initial-data} 所给的形式. 首先将初始曲面化为平面.

\PDEIndexMark{idx:three-surface-descriptions}通常用以下三种方式之一描述曲面:
\begin{enumerate}
\item 作为函数的等值集: $\phi(x_1,\ldots,x_n)=\lambda$. 这里假设 $\nabla\phi\ne0$.
\item 作为图像: $x_i=\psi(x_1,\ldots,x_{i-1},x_{i+1},\ldots,x_n)$.
\item 通过参数化 $\bm{x}=\bm{F}(\bm{y})$, 其中 $\bm{y}\in\mathbb{R}^{n-1}$. 这里假设矩阵 $\partial F_i/\partial y_j$ 具有最大秩.
\end{enumerate}

如果上述三种情形中各自对应的函数 $\phi$, $\psi$, $\bm{F}$ 属于 $C^k$ 类 (解析), 则称曲面属于 $C^k$ 类, $k\geq1$ (解析). 隐函数定理表明, 这三个定义在局部等价. 因而, 必要时重新编号坐标后, 可假设解析曲面由 $x_n=\psi(x_1,\ldots,x_{n-1})$ 给出, 其中 $\psi$ 解析. 现在令 $i\leq n-1$ 时 $y_i=x_i$, 并令 $y_n=x_n-\psi(x_1,\ldots,x_{n-1})$. 在新坐标下, 初始曲面为 $y_n=0$. 因此, 以下假设初始曲面是平面 $x_n=0$.

化简方程时, 有时必须将方程对 $x_n$ 求导. 注意, 这样求导以后, 平面 $x_n=0$ 仍是新方程组的非特征曲面. 此外, 如果求导后的方程处处成立, 且原方程在 $x_n=0$ 时成立, 那么原方程也处处成立. 即每当将方程对 $x_n$ 求导时, 都把原方程看作新方程组的初始数据必须满足的约束.

注意, 将一个或多个方程对 $x_n$ 求导, 就可以使任意偏微分方程组成为拟线性方程组. 因而只需考虑拟线性方程组. 考虑形如
\begin{equation}\label{eq:general-weighted-quasilinear-system}
\sum_{j=1}^N\sum_{|\alpha|=s_i+t_j}a_{ij}^{\alpha}(\bm{p}_i)D^\alpha u_j+b_i(\bm{p}_i)=0,\qquad i=1,\ldots,N,
\end{equation}
的方程组, 其中 $s_i$ 和 $t_j$ 分别是赋给方程与因变量的权 (参见 Example~\ref{ex:weighted-laplace-reformulation} 中的讨论), 向量 $\bm{p}_i$ 由所有自变量以及因变量 $u_j$ 的所有低于 $s_i+t_j$ 阶的导数组成. 由于只有和 $s_i+t_j$ 有关, 可以假设所有 $s_i$ 非正, 所有 $t_j$ 非负. 将第 $i$ 个方程对 $x_n$ 求导 $-s_i$ 次, 就可使新方程组中的 $s_i$ 等于零. 这样做以后, 通过指定 $l=0,\ldots,t_j-1$ 时的 $\partial^l u^j/\partial x_n^l$, 得到一组自然的初始条件. 这些初始条件必须满足对 $x_n$ 求过导的各个方程所产生的约束.

至此得到形如
\begin{equation}\label{eq:zero-row-weight-system}
\sum_{j=1}^N\sum_{|\alpha|=t_j}a_{ij}^{\alpha}(\bm{p})D^\alpha u_j+b_i(\bm{p})=0,\qquad i=1,\ldots,N,
\end{equation}
的方程组, 其中 $\bm{p}$ 包含所有自变量以及所有因变量 $u_j$ 的所有低于 $t_j$ 阶的导数. 初始条件为
\begin{equation}\label{eq:weighted-system-initial-derivatives}
\dfrac{\partial^l u_i}{\partial x_n^l}=f_{il}(x_1,\ldots,x_{n-1}),\qquad i=1,\ldots,N;\quad l=0,\ldots,t_i-1.
\end{equation}


注意, 向量 $\bm{p}$ 由初始数据确定. 如果 (对给定数据而言) 曲面 $x_n=0$ 是非特征的, 就可由 \eqref{eq:zero-row-weight-system} 解出导数
\begin{equation}\label{eq:highest-normal-derivatives}
\dfrac{\partial^{t_i}u_i}{(\partial x_n)^{t_i}},
\end{equation}
并得到形如
\begin{equation}\label{eq:normal-solved-weighted-system}
\dfrac{\partial^{t_i}u_i}{(\partial x_n)^{t_i}}=\sum_{j=1}^N\sum_{\substack{|\alpha|=t_j\\\alpha\ne(0,\ldots,t_j)}}c_{ij}^{\alpha}(\bm{p})D^\alpha u_j+d_i(\bm{p}).
\end{equation}
的方程组.

现在将 $u_j$ 的所有低于 $t_j$ 阶的导数引入为新因变量, 就能将此方程组化为一阶. 然后利用 \eqref{eq:normal-solved-weighted-system} 及混合偏导数相等, 相应地加入新方程. 例如, 若 $t_1=2$, 就引入新变量 $v_{1k}=\partial u_1/\partial x_k$, 并加入方程
\begin{equation}\label{eq:reduction-first-order-derivative-relations}
\dfrac{\partial u_1}{\partial x_n}=v_{1n},\qquad\dfrac{\partial v_{1k}}{\partial x_n}=\dfrac{\partial v_{1n}}{\partial x_k},\qquad k=1,\ldots,n-1,
\end{equation}
而 $\partial v_{1n}/\partial x_n$ 的方程来自 \eqref{eq:normal-solved-weighted-system}.

如此便得到一阶方程组的初值问题. 与 \eqref{eq:ck-intro-quasilinear-system}, \eqref{eq:ck-intro-zero-initial-data} 之间仅有的区别, 是初始数据可能不为零, 且系数可能依赖于 $x_n$. 前者可通过代换处理. 如果 $u_i$ 的初始条件为 $u_i=f_i(x_1,\ldots,x_{n-1})$, 只需引入 $v_i=u_i-\allowbreak f_i(x_1,\ldots,x_{n-1})$ 为新变量. 最后, 引入额外因变量 $w$, 使它满足方程
\begin{equation}\label{eq:reduction-extra-independent-coordinate}
\dfrac{\partial w}{\partial x_n}=1,\qquad w=0\quad\text{on }x_n=0.
\end{equation}
就能消去系数对自变量 $x_n$ 的依赖. 对所有依赖于 $x_n$ 的系数, 用 $w$ 替换 $x_n$.

\begin{Example}\label{ex:scalar-general-reduction}考虑标量方程
\begin{equation}\label{eq:scalar-nonlinear-reduction-origin}
\phi(u_{xx},u_{xy},u_{yy},u_x,u_y,u,x,y)=0
\end{equation}
初始条件为
\begin{equation}\label{eq:scalar-nonlinear-reduction-initial}
u(x,0)=f(x),\qquad u_y(x,0)=g(x).
\end{equation}
假设 $\phi$, $f$ 和 $g$ 都是其自变量的解析函数, 并且至少在原点的一个邻域内, 平面 $y=0$ 是非特征的, 即方程
\begin{equation}\label{eq:scalar-nonlinear-initial-constraint}
\phi(f''(x),g'(x),z,f'(x),g(x),f(x),x,0)=0
\end{equation}
具有解 $z(x)$, 并且
\begin{equation}\label{eq:scalar-nonlinear-noncharacteristic}
\phi_{,3}(f''(x),g'(x),z(x),f'(x),g(x),f(x),x,0)\ne0.
\end{equation}
这里及下文中, $\phi_{,i}$ 表示 $\phi$ 对其第 $i$ 个自变量的导数.

首先对 $y$ 求导, 使方程成为拟线性的. 得到问题
\begin{equation}\label{eq:differentiated-scalar-nonlinear-equation}
\begin{aligned}
0&=\phi_{,1}(\cdots)u_{xxy}+\phi_{,2}(\cdots)u_{xyy}+\phi_{,3}(\cdots)u_{yyy}\\
&\quad+\phi_{,4}(\cdots)u_{xy}+\phi_{,5}(\cdots)u_{yy}\\
&\quad+\phi_{,6}(\cdots)u_y+\phi_{,8}(\cdots),
\end{aligned}
\end{equation}
(其中省略号 $(\cdots)$ 表示自变量 $(u_{xx},u_{xy},u_{yy},u_x,u_y,u,x,y)$), 初始条件为
\begin{equation}\label{eq:differentiated-scalar-initial-data}
u(x,0)=f(x),\quad u_y(x,0)=g(x),\quad u_{yy}(x,0)=z(x),
\end{equation}
这里 $z(x)$ 同上.

现在引入新变量 $p=u_x$, $q=u_y$, $a=u_{xx}$, $b=u_{xy}$, $c=u_{yy}$, $y=w$, 得到方程组
\begin{equation}\label{eq:first-order-scalar-nonlinear-system}
\begin{aligned}
u_y&=q,\quad p_y=b,\quad q_y=c,\quad a_y=b_x,\quad b_y=c_x,\\
c_y&=-\dfrac1{\phi_{,3}(\cdots)}
\bigl[\phi_{,1}(\cdots)b_x+\phi_{,2}(\cdots)c_x\\
&\qquad+\phi_{,4}(\cdots)b+\phi_{,5}(\cdots)c+\phi_{,6}(\cdots)q+\phi_{,8}(\cdots)\bigr],\\
w_y&=1.
\end{aligned}
\end{equation}
这里省略号 $(\cdots)$ 表示自变量 $(a,b,c,p,q,u,x,w)$. 初始条件为
\begin{equation}\label{eq:first-order-scalar-nonlinear-initial-data}
\begin{gathered}
u(x,0)=f(x),\quad p(x,0)=f'(x),\quad q(x,0)=g(x),\quad a(x,0)=f''(x),\\
b(x,0)=g'(x),\quad c(x,0)=z(x),\quad w(x,0)=0.
\end{gathered}
\end{equation}
\end{Example}

\begin{Example}\label{ex:stokes-general-reduction}作为方程组的例子, 考虑二维空间中的 Stokes 方程组
\begin{equation}\label{eq:stokes-two-dimensional-reduction-origin}
\begin{aligned}
u_{xx}+u_{yy}-p_x&=0,\\v_{xx}+v_{yy}-p_y&=0,\\u_x+v_y&=0,
\end{aligned}
\end{equation}
及 $y=0$ 上的初始条件. 约定变量次序为 $(u,v,p)$, 可赋予权 $s_1=s_2=0$, $s_3=-1$, $t_1=t_2=2$, $t_3=1$. 按照上述算法, 必须将 \eqref{eq:stokes-two-dimensional-reduction-origin} 的最后一个方程对 $y$ 求导, 得到
\begin{equation}\label{eq:stokes-differentiated-divergence}
u_{xy}+v_{yy}=0.
\end{equation}
然后将方程组重写为
\begin{equation}\label{eq:stokes-normal-solved-system}
\begin{aligned}
u_{yy}&=-u_{xx}+p_x,\\v_{yy}&=-u_{xy},\\p_y&=v_{xx}-u_{xy}.
\end{aligned}
\end{equation}


适当的初始条件为
\begin{equation}\label{eq:stokes-reduction-initial-data}
\begin{gathered}
u(x,0)=f_1(x),\quad u_y(x,0)=f_2(x),\\
v(x,0)=g_1(x),\quad v_y(x,0)=g_2(x),\\
p(x,0)=h(x);
\end{gathered}
\end{equation}
然而, 为得到原方程组的解, 这些初始数据必须满足 \eqref{eq:stokes-two-dimensional-reduction-origin} 的最后一个方程给出的约束, 即
\begin{equation}\label{eq:stokes-reduction-initial-constraint}
f_1'(x)+g_2(x)=0.
\end{equation}
为得到一阶方程组, 最后令 $u_x=a$, $u_y=b$, $v_x=c$, $v_y=d$. 得到
\begin{equation}\label{eq:stokes-first-order-reduction}
\begin{gathered}
u_y=b,\quad a_y=b_x,\quad b_y=-a_x+p_x,\quad v_y=d,\\
c_y=d_x,\quad d_y=-b_x,\quad p_y=c_x-b_x.
\end{gathered}
\end{equation}
\end{Example}
